% Hybrid Fuzzy-Bayesian Engine

\section{Hybrid Fuzzy-Bayesian Engine}

\subsection{Evidence Fuzzification Process}

Input evidence vector $\mathbf{e} = (e_1, e_2, \ldots, e_n)^T$ undergoes fuzzification using triangular membership functions \cite{dubois1980fuzzy}. For evidence component $e_i \in [0, 1]$:

Low confidence membership function:
\begin{equation}
\mu_{LOW}(e_i) = \begin{cases}
1 & \text{if } e_i \leq 0.3 \\
\frac{0.5 - e_i}{0.5 - 0.3} & \text{if } 0.3 < e_i < 0.5 \\
0 & \text{if } e_i \geq 0.5
\end{cases}
\end{equation}

Medium confidence membership function:
\begin{equation}
\mu_{MEDIUM}(e_i) = \begin{cases}
0 & \text{if } e_i \leq 0.3 \\
\frac{e_i - 0.3}{0.5 - 0.3} & \text{if } 0.3 < e_i < 0.5 \\
1 & \text{if } e_i = 0.5 \\
\frac{0.7 - e_i}{0.7 - 0.5} & \text{if } 0.5 < e_i < 0.7 \\
0 & \text{if } e_i \geq 0.7
\end{cases}
\end{equation}

High confidence membership function:
\begin{equation}
\mu_{HIGH}(e_i) = \begin{cases}
0 & \text{if } e_i \leq 0.5 \\
\frac{e_i - 0.5}{0.7 - 0.5} & \text{if } 0.5 < e_i < 0.7 \\
1 & \text{if } e_i \geq 0.7
\end{cases}
\end{equation}

\subsection{Fuzzy Rule Base}

Molecular identification rules using Mamdani fuzzy inference \cite{mamdani1975experiment}. Rule structure:
\begin{equation}
R_k: \text{IF } e_1 \text{ is } A_{1k} \text{ AND } e_2 \text{ is } A_{2k} \text{ THEN } M \text{ is } B_k
\end{equation}

where $A_{ik} \in \{LOW, MEDIUM, HIGH\}$ and $B_k$ represents molecular identity confidence.

Rule activation strength using minimum t-norm:
\begin{equation}
\alpha_k = \min(\mu_{A_{1k}}(e_1), \mu_{A_{2k}}(e_2), \ldots, \mu_{A_{nk}}(e_n))
\end{equation}

Output fuzzy set for rule $k$:
\begin{equation}
\mu_{B_k}^{output}(y) = \min(\alpha_k, \mu_{B_k}(y))
\end{equation}

\subsection{Bayesian Network Integration}

Prior probability distribution for molecular identities \cite{laplace1814essai}:
\begin{equation}
P(M = m_j) = \frac{N_j}{\sum_{i=1}^K N_i}
\end{equation}

where $N_j$ is frequency count for molecule $m_j$ in training database and $K$ is total number of molecular identities.

Likelihood calculation incorporating fuzzy evidence:
\begin{equation}
P(E = e_i | M = m_j) = \int_0^1 P(E = e_i | M = m_j, \mu = u) \cdot P(\mu = u | M = m_j) \, du
\end{equation}

where $\mu$ represents membership degree and integration accounts for fuzzy uncertainty.

Simplified likelihood using centroid defuzzification:
\begin{equation}
P(E = e_i | M = m_j) = \mathcal{N}(e_i; \mu_{ij}, \sigma_{ij}^2)
\end{equation}

where $\mathcal{N}$ denotes normal distribution with mean $\mu_{ij}$ and variance $\sigma_{ij}^2$ learned from training data.

\subsection{Posterior Probability Calculation}

Fuzzy-Bayesian posterior using weighted evidence combination:
\begin{equation}
P(M = m_j | \mathbf{e}) = \frac{\sum_{l=1}^L w_l \cdot P(\mathbf{e}_l | M = m_j) \cdot P(M = m_j)}{\sum_{k=1}^K \sum_{l=1}^L w_l \cdot P(\mathbf{e}_l | M = m_k) \cdot P(M = m_k)}
\end{equation}

where:
- $L$ is number of fuzzy evidence categories
- $w_l$ are fuzzy membership weights
- $\mathbf{e}_l$ are evidence vectors corresponding to fuzzy categories

Weight calculation using normalized membership degrees:
\begin{equation}
w_l = \frac{\mu_l}{\sum_{i=1}^L \mu_i}
\end{equation}

\subsection{Temporal Evidence Integration}

Time-dependent evidence reliability using exponential decay \cite{cox1962renewal}:
\begin{equation}
R_i(t) = R_{i0} \cdot \exp(-\lambda_i t)
\end{equation}

where:
- $R_{i0}$ is initial reliability for evidence source $i$
- $\lambda_i$ is decay constant for evidence type $i$
- $t$ is time elapsed since evidence collection

Time-adjusted posterior probability:
\begin{equation}
P(M = m_j | \mathbf{e}, \mathbf{t}) = \frac{\sum_{l=1}^L w_l \cdot R_l(t_l) \cdot P(\mathbf{e}_l | M = m_j) \cdot P(M = m_j)}{\sum_{k=1}^K \sum_{l=1}^L w_l \cdot R_l(t_l) \cdot P(\mathbf{e}_l | M = m_k) \cdot P(M = m_k)}
\end{equation}

\subsection{Network Coherence Optimization}

Evidence network represented as graph $G = (V, E)$ where vertices $V$ represent evidence nodes and edges $E$ represent relationships \cite{newman2003structure}. Coherence measure:
\begin{equation}
\mathcal{H} = \frac{1}{|E|} \sum_{(i,j) \in E} \cos(\theta_{ij})
\end{equation}

where $\theta_{ij}$ is angle between evidence vectors $\mathbf{e}_i$ and $\mathbf{e}_j$:
\begin{equation}
\cos(\theta_{ij}) = \frac{\mathbf{e}_i \cdot \mathbf{e}_j}{||\mathbf{e}_i|| \cdot ||\mathbf{e}_j||}
\end{equation}

Coherence optimization objective function:
\begin{equation}
\max_{\mathbf{w}} \mathcal{H}(\mathbf{w}) \text{ subject to } \sum_{i=1}^L w_i = 1, \quad w_i \geq 0
\end{equation}

Solution using Lagrangian method:
\begin{equation}
L = \mathcal{H}(\mathbf{w}) - \lambda \left(\sum_{i=1}^L w_i - 1\right) - \sum_{i=1}^L \mu_i w_i
\end{equation}

Optimal weights satisfy:
\begin{equation}
\frac{\partial \mathcal{H}}{\partial w_i} = \lambda + \mu_i
\end{equation}

\subsection{Uncertainty Quantification}

Evidence uncertainty propagation through fuzzy-Bayesian network \cite{klir1995fuzzy}. Input uncertainty for evidence component $i$:
\begin{equation}
\sigma_{input,i}^2 = \sigma_{measurement,i}^2 + \sigma_{model,i}^2
\end{equation}

Fuzzy membership uncertainty:
\begin{equation}
\sigma_{\mu,i}^2 = \left(\frac{\partial \mu}{\partial e_i}\right)^2 \sigma_{input,i}^2
\end{equation}

Posterior probability uncertainty using error propagation:
\begin{equation}
\sigma_{posterior}^2 = \sum_{i=1}^L \left(\frac{\partial P(M|\mathbf{e})}{\partial w_i}\right)^2 \sigma_{w,i}^2
\end{equation}

where $\sigma_{w,i}^2$ represents uncertainty in fuzzy weight $w_i$.

Confidence interval for molecular identification:
\begin{equation}
CI_{1-\alpha} = P(M = m_j | \mathbf{e}) \pm z_{\alpha/2} \cdot \sigma_{posterior}
\end{equation}

\subsection{Defuzzification and Decision Making}

Centroid defuzzification for final molecular identification:
\begin{equation}
m^* = \frac{\sum_{j=1}^K m_j \cdot P(M = m_j | \mathbf{e})}{\sum_{j=1}^K P(M = m_j | \mathbf{e})}
\end{equation}

Decision threshold based on maximum posterior probability:
\begin{equation}
\text{Decision} = \begin{cases}
m^* & \text{if } \max_j P(M = m_j | \mathbf{e}) > \theta_{decision} \\
\text{Unknown} & \text{otherwise}
\end{cases}
\end{equation}

where $\theta_{decision} = 0.5$ is minimum confidence threshold for positive identification.
